\chapter{Wiener--Laguerre scalar criterion}\label{ch:wiener-laguerre-scalar}

\section{Reduction of the density frontier}

Chapter~\ref{ch:correlation-kernel} expresses the Dimitrov--Xu critical-axis criterion as density of the translates of
\[
\Phi_{2,y}(t)=\cosh(ty)\nu_2(t),
\qquad
0<|y|<\frac12.
\]
Wiener's $L^1$ translation theorem allows this infinite-dimensional statement to be reduced to a pointwise condition on a Fourier transform.  The Wronskian identities of Dimitrov and Xu then identify that transform directly with derivatives of the completed Xi function \cite{DimitrovXu2016}.

For real $x$ and admissible $y$, define
\begin{equation}\label{eq:q-xi-definition}
Q_\Xi(x,y)
:=
|\Xi'(x+iy)|^2
-\Ree\!\left(
\Xi(x+iy)\overline{\Xi''(x+iy)}
\right).
\end{equation}

\section{Wronskian identity}

Write
\[
\Xi(x+iy)=\phi(x,y)+i\psi(x,y).
\]
The Cauchy--Riemann equations give
\begin{align}
Q_\Xi(x,y)
={}&\phi_x(x,y)^2-\phi(x,y)\phi_{xx}(x,y)\\
&+\psi_x(x,y)^2-\psi(x,y)\psi_{xx}(x,y).
\end{align}
With $W_2(f)=ff''-(f')^2$, this becomes
\begin{equation}\label{eq:q-xi-wronskian}
Q_\Xi(x,y)
=-W_2(\phi(\cdot,y);x)-W_2(\psi(\cdot,y);x).
\end{equation}
Dimitrov--Xu Theorem~1.3 / Theorem~2.5 and their Section~3 correlation identity yield
\begin{equation}\label{eq:q-xi-fourier}
\boxed{
Q_\Xi(x,y)=\mathcal F[\Phi_{2,y}](x).
}
\end{equation}
The repository records Eq.~\eqref{eq:q-xi-fourier} as \status{STANDARD\_PASS}.

\section{Tauberian equivalence}

For each fixed admissible $y$, Wiener's $L^1$ theorem states
\begin{equation}\label{eq:wiener-density-nonzero}
\overline{\operatorname{span}\mathcal T(\Phi_{2,y})}^{L^1(\mathbb R)}
=L^1(\mathbb R)
\iff
\mathcal F[\Phi_{2,y}](x)\neq0
\quad\forall x\in\mathbb R.
\end{equation}
The Dimitrov--Xu proof also gives strict positivity at the origin,
\begin{equation}\label{eq:q-origin-positive}
Q_\Xi(0,y)>0
\qquad
\left(0<|y|<\frac12\right).
\end{equation}
Because $Q_\Xi(\cdot,y)$ is continuous and real-valued on $\mathbb R$, a globally nonvanishing $Q_\Xi$ has the positive sign selected by Eq.~\eqref{eq:q-origin-positive}.  Therefore
\begin{equation}\label{eq:density-q-positive}
\boxed{
\overline{\operatorname{span}\mathcal T(\Phi_{2,y})}^{L^1}=L^1
\iff
Q_\Xi(x,y)>0\quad\forall x\in\mathbb R.
}
\end{equation}
Combining Eq.~\eqref{eq:density-q-positive} with Dimitrov--Xu Theorem~1.1 gives
\begin{theorem}[Wiener--Laguerre scalar RH criterion]\label{thm:wiener-laguerre-rh}
Under the standard Xi-kernel hypotheses used by Dimitrov and Xu,
\begin{equation}\label{eq:q-rh-equivalence}
\boxed{
\mathrm{RH}
\iff
\forall\,0<|y|<\frac12\ \forall x\in\mathbb R:\quad
Q_\Xi(x,y)>0.
}
\end{equation}
\end{theorem}

The equivalence in Theorem~\ref{thm:wiener-laguerre-rh} has status \status{STANDARD}.  The globally quantified sign condition has status
\[
\status{OPEN\_RH\_EQUIVALENT\_CRITERION}.
\]

\section{Executable diagnostic}

The function
\begin{center}
\texttt{xi\_wiener\_laguerre\_scalar(x, y)}
\end{center}
implements Eq.~\eqref{eq:q-xi-definition} and enforces the open domain $0<|y|<1/2$.  The tests verify $y\leftrightarrow-y$ symmetry, positive origin reference values, and fail-closed domain handling.

Finite samples retain \status{NUMERICAL\_DIAGNOSTIC} authority.  An analytic closure must carry both universal quantifiers in Eq.~\eqref{eq:q-rh-equivalence}.

\section{Certified falsification surface}

The scalar form provides a direct falsification target.  A certified pair $(x_*,y_*)$ with
\[
x_*\in\mathbb R,
\qquad
0<|y_*|<\frac12,
\qquad
Q_\Xi(x_*,y_*)\le0
\]
combined with Eq.~\eqref{eq:q-origin-positive} and continuity would force a real zero of $Q_\Xi(\cdot,y_*)=\mathcal F[\Phi_{2,y_*}]$.  The corresponding Wiener density condition would fail at that $y_*$.

This makes the next research frontier sharply local in representation while globally quantified in domain:
\begin{equation}\label{eq:xf4-frontier}
\boxed{
Q_\Xi(x,y)>0
\quad
\text{for every }x\in\mathbb R,
\quad0<|y|<\frac12.
}
\end{equation}
The next analytic task is to seek a structural representation of $Q_\Xi$ that certifies this sign over the entire strip.
